% ==============================================================================
%  Conviction-based inference — supplementary material
%
%  Everything here is reproducible from the committed per-dataset tables under
%  quick_analysis/Conviction_experiments_tables/ via
%      python quick_analysis/conviction_summary_tables.py \
%             --score {last|avg|avg_z} --exclude cornell,wisconsin \
%             --out_dir quick_analysis/conviction_tables
%  which writes one .tex file per (experiment, intervention), each already
%  carrying its own caption and \label, plus the matching CSVs.
%
%  Path convention follows the Forman section: quick_analysis/conviction_tables/
%  is sma_projet_files/Sections/conviction_tables/, and
%  quick_analysis/Norm_Analysis/ is Deck_slides/.
% ==============================================================================

\section{Supplementary material: conviction-based inference}
\label{sec:conviction_supp}

This appendix collects the material that supports
Section~\ref{sec:conviction} but would obstruct its argument: the scope of the
equilibrium prediction, the per-dataset figures behind the observational claims,
the complete tables for the three perturbation modes not shown in the main text,
the robustness check under a different conviction score, the implementation
details of each intervention, and the diagnostics that certify that the pruned
graphs remain usable objects.

% ------------------------------------------------------------------------------
\subsection{Scope of the equilibrium prediction}
\label{sec:supp_conv_scope}

The equilibrium result that Section~\ref{sec:conviction} tests is proved in a
setting slightly narrower than the one we measure, and we record the discrepancy
here rather than in the main text. In the theorem the state of a node is an
element of $\mathbb{R}^d$, whereas the hidden state of a Sheaf Neural Network is
a matrix $X_t[u] \in \mathbb{R}^{d \times F}$, since the network carries $F$
feature channels alongside the $d$ stalk dimensions and interleaves the
diffusion with weight matrices. Our setting is therefore richer than the one in
which the equilibrium is derived, and the prediction should be read as a guide
rather than as a statement that applies verbatim.

We nonetheless regard the comparison as meaningful, for two reasons. The
quantity the theorem constrains, the magnitude of a node's representation, is
exactly the one we measure; and the conclusion we draw from the comparison is
negative, namely that the process does not do what the equilibrium prescribes.
A negative conclusion is the direction in which the mismatch is harmless: a
richer state space gives the process more room to depart from the prediction,
not less, so observing that it departs is not an artefact of the extra
dimensions. What we could not do on this evidence is the converse, and we make
no claim that would require the theorem to hold verbatim.

% ------------------------------------------------------------------------------
\subsection{What the conviction correlates with}
\label{sec:supp_conv_observational}

Section~\ref{sec:conviction} settles by intervention the question of whether the
conviction identifies the nodes a trained network relies on, and we stress there
that no correlation could have settled it. It is nonetheless worth recording
what the conviction does correlate with, both because it is the natural first
question a reader will ask and because the untrained checkpoint tells us
something the trained one cannot. Table~\ref{tab:conv_observational} collects
the six observables, averaged over the folds of a dataset and then over the
eight datasets, at the last diffusion layer.

The epoch-$0$ column of Table~\ref{tab:conv_observational} is what makes each
configuration legible, since it says how
much of what we measure is already present before any learning has happened. In
the normalised general sheaf, the configuration the main text is about, the
correlation with degree is absent at initialisation and builds up gradually,
growing with depth and with the number of epochs, so whatever ordering the
conviction ends up expressing is produced by training rather than handed down by
the architecture. We would be accurate about its size: the conviction of the
general sheaf is not \emph{uncorrelated} with degree, it is moderately
correlated with it. Roman Empire is a representative case, going from $-0.143$
to $0.432$ across the layers (Figure~\ref{fig: RM_General_Scatter}). Of the
other observables, the conviction is most strongly associated with the model's
softmax confidence, at $0.66 \pm 0.05$, more strongly than with either
topological quantity, while its correlation with input-gradient saliency is
modest and, alone among the six, negative on average. The Cohen's $d$ separating correctly from
incorrectly classified nodes is moderate at the best epoch; its trajectory along
the checkpoint axis is not the same on every dataset, growing with training on
some and peaking in the first epochs on others, so we do not read a common
pattern into it.

In the unnormalised bundle sheaf the conviction is instead coupled to the degree
from the start. Reading the second column group of
Table~\ref{tab:conv_observational}, the coupling is already $0.57 \pm 0.08$ in the untrained
network, and training barely moves it, to $0.53 \pm 0.09$; a single diffusion
step is enough to produce it, as Figure~\ref{fig: Tolokers_Bundle_Scatter} shows
for the most extreme dataset, where the correlation is $0.956$ at $X_1$. This is
the behaviour one should expect of an unnormalised Laplacian, whose operator carries
the degree in its own scale, and it is consistent with the reading of
Section~\ref{sec:conviction}: the norms are ordered by the topology rather than
by the task, which is a further way of saying that the process has not settled
where the equilibrium result would put it.

We would outline that the two configurations differ both in the family of
restriction maps and in the normalisation, so nothing in
Table~\ref{tab:conv_observational} isolates either factor, and we do not read
the difference between its column groups as an ablation. What the table gives is
a description of what each pipeline does.

% Produced by:
%   python quick_analysis/conviction_heatmap.py --model GeneralSheaf --normalised true
%   python quick_analysis/conviction_heatmap.py --model BundleSheaf  --normalised false
%   python quick_analysis/conviction_observational_table.py
% or, in one job:  sbatch exp/scripts/run_conviction_heatmap_epfl.sh
% auto-generated by quick_analysis/conviction_summary_tables.py
\begin{table}[H]
\centering
\small
\begin{tabular}{lcccc}
\toprule
 & \multicolumn{2}{c}{General sheaf, $\widetilde{\Delta}_{\mathcal{F}}$} & \multicolumn{2}{c}{Bundle sheaf, $\Delta_{\mathcal{F}}$} \\
\cmidrule(lr){2-3}\cmidrule(lr){4-5}
Observable & Untrained & Best epoch & Untrained & Best epoch \\
\midrule
Spearman $\rho$ with degree & $0.07 \pm 0.10$ & $0.33 \pm 0.07$ & $0.57 \pm 0.08$ & $0.53 \pm 0.09$ \\
Spearman $\rho$ with betweenness & $0.05 \pm 0.08$ & $0.28 \pm 0.06$ & $0.46 \pm 0.07$ & $0.45 \pm 0.07$ \\
Spearman $\rho$ with confidence & $0.17 \pm 0.07$ & $0.66 \pm 0.05$ & $0.53 \pm 0.13$ & $0.55 \pm 0.07$ \\
Spearman $\rho$ with saliency & $0.07 \pm 0.10$ & $-0.18 \pm 0.09$ & $0.82 \pm 0.10$ & $0.65 \pm 0.12$ \\
$|$Cohen's $d|$, train split & $0.21 \pm 0.05$ & $0.54 \pm 0.09$ & $0.10 \pm 0.04$ & $0.17 \pm 0.05$ \\
$|$Cohen's $d|$, test split & $0.24 \pm 0.05$ & $0.37 \pm 0.08$ & $0.16 \pm 0.06$ & $0.20 \pm 0.05$ \\
\bottomrule
\end{tabular}
\caption{Observational profile of the conviction at the last diffusion layer $X_L$. Each entry is a Spearman rank correlation between the conviction and the named node property, or the absolute Cohen's $d$ separating the conviction of correctly and incorrectly classified nodes, computed per fold, averaged over the folds of a dataset, and then averaged over the 8 datasets. Mean $\pm$ standard error of the cross-dataset mean, the quoted spread being between-dataset variability. ``Untrained'' is the epoch-$0$ checkpoint, where the sheaf is randomly initialised and nothing has been learnt, and it therefore acts as a control: an entry that is already large there is produced by the architecture rather than by training. The two configurations differ in the restriction-map family and in the normalisation at once, so the table describes what each pipeline does and is not an ablation of either factor.}
\label{tab:conv_observational}
\end{table}


% ------------------------------------------------------------------------------
\subsection{Per-dataset observational figures}
\label{sec:supp_conv_scatter}

Figure~\ref{fig: RM_General_Scatter} is the counterpart, for the normalised
general sheaf, of Figure~\ref{fig: RM_Bundle_Scatter} of
Section~\ref{sec:conviction}. Comparing the two on the same dataset is the
cleanest way of seeing the difference discussed in the main text: both models
start from essentially no relation between conviction and degree on the input
features, and in both the relation strengthens with each diffusion step, but the
bundle sheaf reaches a markedly higher correlation than the general sheaf does.

\begin{figure}[H]
    \centering
    \includegraphics[width=\linewidth]{Deck_slides/conviction_scatter/GeneralSheaf/normalised_true/roman_empire_GeneralSheaf_d5_5L_Norm_True_conviction_scatter.pdf}
    \caption{Conviction against node degree, one panel per diffusion layer, for
    the normalised general sheaf on Roman Empire ($d=5$, $L=5$). Norms are taken
    at the best-validation-epoch checkpoint and averaged over the ten folds; the
    dashed line is a linear fit and the annotated $\rho_s$ is the Spearman rank
    correlation with degree. The correlation grows from $-0.143$ at $X_0$ to
    $0.432$ at $X_5$, against $-0.116$ and $0.652$ for the bundle sheaf on the
    same graph (Figure~\ref{fig: RM_Bundle_Scatter}).}
    \label{fig: RM_General_Scatter}
\end{figure}

Figure~\ref{fig: Tolokers_Bundle_Scatter} shows the most extreme case of the
bundle-sheaf behaviour in our suite. On Tolokers the conviction of the
unnormalised bundle sheaf is almost a monotone function of the degree from the
very first diffusion step onwards, which is the clearest illustration of the
point made in the main text: in an unnormalised Laplacian the degree enters the
scale of the operator itself, so the ordering of the stalk norms is fixed by the
topology rather than learnt from the task.

\begin{figure}[H]
    \centering
    \includegraphics[width=\linewidth]{Deck_slides/conviction_scatter/BundleSheaf/normalised_false/tolokers_BundleSheaf_d3_4L_Norm_False_conviction_scatter.pdf}
    \caption{Conviction against node degree for the unnormalised bundle sheaf on
    Tolokers ($d=3$, $L=4$), at the best-validation-epoch checkpoint and averaged
    over folds. The Spearman correlation with degree is $-0.119$ on the input
    features and $0.956$ already after one diffusion step, settling at $0.881$ at
    the last layer.}
    \label{fig: Tolokers_Bundle_Scatter}
\end{figure}

Both the figures above and Table~\ref{tab:conv_observational} come from the same
computation, which we describe here in full since the main text only summarises
its outcome. For each dataset and each configuration we compare the conviction
against other node properties at every training checkpoint and every diffusion
layer, producing a grid of heatmaps whose rows are checkpoints and whose columns
are layers. Each cell reports one scalar: the Spearman correlation of
$\mathbf{c}_t$ with node degree, with betweenness centrality, with the model's
softmax confidence and with input-gradient saliency, together with the absolute
Cohen's $d$ separating the conviction of correctly and incorrectly classified
nodes, computed separately on train and test. We take the absolute value of
Cohen's $d$ before averaging over folds, so that sign cancellations across folds
do not hide an effect that is present in each of them; it is undefined when the
training accuracy reaches $100\%$, since the wrong group is then empty, and those
cells are annotated as such rather than imputed. As in the Forman analysis, the
epoch-$0$ row is an untrained sheaf and therefore acts as a control: it tells us
which part of what we measure is produced by the architecture and which part is
produced by learning.

The heatmaps average the folds away before anything reaches disk, so
Table~\ref{tab:conv_observational} is built from the per-fold values written
alongside them, and aggregated exactly as the intervention tables are: per fold,
then averaged over the folds of a dataset, then averaged over datasets with the
spread taken between datasets. The scatter figures are a different projection of
the same quantity: they fix the best-epoch checkpoint, average the node norms
over folds, and plot them against degree one panel per diffusion layer, which is
the most direct way of inspecting whether high-conviction nodes are simply hubs.
The Spearman value annotated on a scatter panel is therefore computed on
fold-averaged norms, whereas the table averages per-fold correlations; the two
answer slightly different questions and need not agree to the last decimal.

% ------------------------------------------------------------------------------
\subsection{Implementation of the interventions}
\label{sec:supp_impl}

All three interventions operate on a model whose weights are loaded from the
best-validation-epoch checkpoint and never updated. The candidate set is always
a subset of the training nodes, and the evaluation set is always the full test
set, so that the quantity being compared is identical across strategies,
budgets and perturbation modes.

\paragraph{Masking.}
The features of the selected nodes are replaced, either by zero or by the
dataset feature mean, and the graph is left untouched. The mean is computed over
all nodes of the graph; this is an unsupervised statistic and involves no label
information, but we note it explicitly since it does touch test-node features.
Because the candidate pool is restricted to training nodes, no test node is ever
masked, and the exclusion of masked nodes from the evaluation index is therefore
a defensive no-op rather than a load-bearing step.

\paragraph{Gating.}
A per-node multiplier $g \in \mathbb{R}^{n}$ is applied to the sparse diffusion
Laplacian, scaling every cross-node entry by the gate of the \emph{source} node.
Entries belonging to a diagonal block, that is those for which row and column
belong to the same node, are left at $1$. A gated node therefore continues to
retain its own state, and only what it transmits to its neighbours is
suppressed. In the hard variant used throughout the main text the gate is $0$
for the selected nodes and $1$ elsewhere. In the soft variant the gate varies
continuously,
\[
    g(u) \;=\; g_{\min} + (1 - g_{\min}) \,
    \frac{s(u) - \min_{v} s(v)}{\max_{v} s(v) - \min_{v} s(v)},
    \qquad g_{\min} = 0.1,
\]
where $s$ is the ranking score, and it is applied to the whole training pool at
once. Passing no gate at all is a complete no-op, so the baseline forward pass
is bit-identical to an ordinary evaluation.

\paragraph{Pruning.}
The selected nodes and their incident edges are deleted and the remaining nodes
are relabelled to a contiguous index space. Since the number of nodes and the
edge index are fixed at model construction, a new model is instantiated for the
reduced graph and the \emph{same} frozen parameters are loaded into it; no
parameter is re-estimated. An assertion guarantees that no validation or test
node is ever among the removed set. This mode is by far the most expensive,
because a model must be rebuilt for every evaluated point.

\paragraph{Degree-matched sampling.}
Given the reference set $R$ (the top- or bottom-$k$ by conviction), the control
set is built by greedy nearest-degree matching without replacement: the
candidate pool is $\mathcal{P} \setminus R$, bucketed by degree and shuffled
once; each $u \in R$ then claims an unused node of the same degree, or of the
closest available degree if that bucket is exhausted. This requires
$k \le |\mathcal{P}|/2$, which every budget we use satisfies; the implementation
raises rather than silently degrading if a larger budget is requested.

\paragraph{Tie-breaking and reproducibility.}
Every ranking is broken by a per-node uniform key drawn from a generator seeded
by the training seed and the fold index. Without this, \texttt{argsort} resolves
ties by ascending node index, which systematically favours low-index nodes at
the boundary of a top-$k$ selection; degree is integer-valued and heavily tied,
so this is not a hypothetical concern. The randomised strategies draw from a
generator keyed on (dataset, fold, strategy, budget) rather than from a single
running stream, so a given draw is reproducible irrespective of which other
datasets or budgets were run before it in the same job.

% ------------------------------------------------------------------------------
\subsection{Aggregation protocol}
\label{sec:supp_agg}

Each per-dataset experiment writes both an aggregated table and a
\texttt{\_folds} table, the latter holding one value per fold. The relative drop
of Section~\ref{sec:conviction} is formed fold by fold, each fold against its
\emph{own} unperturbed baseline, and only then averaged: this keeps the baseline
and the perturbed metric paired within a fold, which is how they were actually
measured. The normalisation is by the headroom above chance,
\[
    \mathrm{drop}\%(s) = 100 \cdot
    \frac{m_{\text{baseline}} - m_{s}}{m_{\text{baseline}} - m_{\text{chance}}},
    \qquad
    m_{\text{chance}} = \begin{cases} 0 & \text{accuracy} \\ 0.5 & \text{ROC-AUC},\end{cases}
\]
so that a cell of any table means the same thing on an accuracy dataset and on
an AUC one, namely the percentage of the learnt signal that the intervention
destroys. Without the floor, an AUC drop divided by the raw baseline would be
understated by roughly a factor of three, since a model at AUC $0.5$ has already
learnt nothing and the raw baseline treats $0$ as the worst attainable score.

Folds are then averaged within a dataset, and the datasets are averaged
unweighted, so each benchmark contributes equally regardless of its size. The
reported dispersion is
\[
    \mathrm{sd} = \sqrt{\tfrac{1}{N-1}\textstyle\sum_i (d_i - \bar d)^2},
    \qquad
    \mathrm{se} = \mathrm{sd}/\sqrt{N},
\]
over the $N$ per-dataset drops $d_i$; it is a \emph{between-dataset} quantity and
deliberately not a fold-level error bar. The ten folds of one dataset share a
graph, an architecture and overlapping splits, so pooling all folds of all
datasets and dividing by $\sqrt{N \cdot 10}$ would produce an error bar several
times too small. The summary CSVs record that naive pooled quantity as
\texttt{total\_se} alongside the reported one, so that the size of the
understatement can be inspected rather than merely asserted.

% ------------------------------------------------------------------------------
\subsection{Complete forward-sweep tables}
\label{sec:supp_forward}

Tables~\ref{tab:conv_fwd_masking-mean_last}--\ref{tab:conv_fwd_pruning_last}
give the budget sweep for the three perturbation modes not shown in the main
text; masking with zero replacement appears as
Table~\ref{tab:conv_fwd_masking-zero_last} there. All four are computed with
conviction score \texttt{last} over the eight retained datasets, and in each of
them the shading marks the most and the least damaging selection at each budget.
High conviction is the most damaging entry in all four tables and at all five
budgets, and low conviction the least, with the single exception of the $1\%$
gating budget, where low degree comes out marginally below it.

% auto-generated by quick_analysis/conviction_summary_tables.py
\begin{table}[H]
\centering
\small
\definecolor{convworst}{rgb}{0.97,0.80,0.78}
\definecolor{convbest}{rgb}{0.80,0.87,0.96}
\begin{tabular}{lccccc}
\toprule
Selection & 1\% & 5\% & 10\% & 20\% & 30\% \\
\midrule
High conviction & \cellcolor{convworst}$0.50 \pm 0.16$ & \cellcolor{convworst}$1.34 \pm 0.39$ & \cellcolor{convworst}$2.19 \pm 0.72$ & \cellcolor{convworst}$3.61 \pm 1.30$ & \cellcolor{convworst}$4.66 \pm 1.85$ \\
Low conviction & \cellcolor{convbest}$0.01 \pm 0.02$ & \cellcolor{convbest}$0.08 \pm 0.05$ & \cellcolor{convbest}$0.14 \pm 0.10$ & \cellcolor{convbest}$0.36 \pm 0.17$ & \cellcolor{convbest}$0.57 \pm 0.38$ \\
High degree & $0.38 \pm 0.16$ & $0.94 \pm 0.38$ & $1.35 \pm 0.59$ & $2.13 \pm 0.94$ & $2.91 \pm 1.17$ \\
Low degree & $0.03 \pm 0.02$ & $0.18 \pm 0.11$ & $0.41 \pm 0.22$ & $0.88 \pm 0.48$ & $1.43 \pm 0.80$ \\
Degree-matched (high) & $0.09 \pm 0.07$ & $0.30 \pm 0.21$ & $0.52 \pm 0.37$ & $0.77 \pm 0.63$ & $1.21 \pm 0.71$ \\
Degree-matched (low) & $0.06 \pm 0.04$ & $0.26 \pm 0.15$ & $0.50 \pm 0.30$ & $1.28 \pm 0.69$ & $2.23 \pm 1.12$ \\
High confidence & $0.15 \pm 0.07$ & $0.51 \pm 0.16$ & $0.86 \pm 0.24$ & $1.57 \pm 0.48$ & $2.25 \pm 0.78$ \\
Low confidence & $0.06 \pm 0.04$ & $0.30 \pm 0.15$ & $0.61 \pm 0.28$ & $1.29 \pm 0.58$ & $1.99 \pm 0.94$ \\
Random & $0.08 \pm 0.04$ & $0.25 \pm 0.16$ & $0.46 \pm 0.32$ & $1.18 \pm 0.60$ & $1.89 \pm 0.94$ \\
\bottomrule
\end{tabular}
\caption{Relative drop (\%) in the headline metric, accuracy or ROC-AUC on Minesweeper, Tolokers and Questions, expressed as a percentage of the baseline's headroom above chance, so that the two metrics are comparable. Intervention: masking (mean), as a function of the fraction of the training pool perturbed. Conviction score \texttt{last}. Each of the 8 datasets contributes one value, itself averaged over its folds. Mean $\pm$ standard error of the cross-dataset mean, the quoted spread being between-dataset variability. Larger is more damaging. In each column the most damaging selection is shaded red and the least damaging blue. }
\label{tab:conv_fwd_masking-mean_last}
\end{table}


% auto-generated by quick_analysis/conviction_summary_tables.py
\begin{table}[H]
\centering
\small
\definecolor{convworst}{rgb}{0.97,0.80,0.78}
\definecolor{convbest}{rgb}{0.80,0.87,0.96}
\begin{tabular}{lccccc}
\toprule
Selection & 1\% & 5\% & 10\% & 20\% & 30\% \\
\midrule
High conviction & \cellcolor{convworst}$0.81 \pm 0.29$ & \cellcolor{convworst}$1.80 \pm 0.41$ & \cellcolor{convworst}$2.77 \pm 0.66$ & \cellcolor{convworst}$4.32 \pm 1.27$ & \cellcolor{convworst}$5.67 \pm 1.89$ \\
Low conviction & $0.05 \pm 0.03$ & \cellcolor{convbest}$0.11 \pm 0.04$ & \cellcolor{convbest}$0.18 \pm 0.09$ & \cellcolor{convbest}$0.43 \pm 0.22$ & \cellcolor{convbest}$0.71 \pm 0.45$ \\
High degree & $0.65 \pm 0.20$ & $1.37 \pm 0.37$ & $1.91 \pm 0.45$ & $2.83 \pm 0.73$ & $3.67 \pm 1.04$ \\
Low degree & \cellcolor{convbest}$0.05 \pm 0.02$ & $0.26 \pm 0.12$ & $0.58 \pm 0.27$ & $1.20 \pm 0.59$ & $1.91 \pm 0.98$ \\
Degree-matched (high) & $0.12 \pm 0.05$ & $0.37 \pm 0.19$ & $0.68 \pm 0.35$ & $1.02 \pm 0.60$ & $1.44 \pm 0.86$ \\
Degree-matched (low) & $0.06 \pm 0.03$ & $0.31 \pm 0.18$ & $0.62 \pm 0.38$ & $1.53 \pm 0.84$ & $2.71 \pm 1.36$ \\
High confidence & $0.24 \pm 0.08$ & $0.75 \pm 0.20$ & $1.22 \pm 0.29$ & $2.28 \pm 0.59$ & $3.27 \pm 0.98$ \\
Low confidence & $0.09 \pm 0.05$ & $0.35 \pm 0.17$ & $0.63 \pm 0.29$ & $1.29 \pm 0.60$ & $2.03 \pm 0.95$ \\
Random & $0.09 \pm 0.04$ & $0.34 \pm 0.17$ & $0.63 \pm 0.34$ & $1.46 \pm 0.69$ & $2.21 \pm 1.03$ \\
\bottomrule
\end{tabular}
\caption{Relative drop (\%) in the headline metric, accuracy or ROC-AUC on Minesweeper, Tolokers and Questions, expressed as a percentage of the baseline's headroom above chance, so that the two metrics are comparable. Intervention: gating (hard), as a function of the fraction of the training pool perturbed. Conviction score \texttt{last}. Each of the 8 datasets contributes one value, itself averaged over its folds. Mean $\pm$ standard error of the cross-dataset mean, the quoted spread being between-dataset variability. Larger is more damaging. In each column the most damaging selection is shaded red and the least damaging blue. }
\label{tab:conv_fwd_gating-hard_last}
\end{table}


% auto-generated by quick_analysis/conviction_summary_tables.py
\begin{table}[H]
\centering
\small
\definecolor{convworst}{rgb}{0.97,0.80,0.78}
\definecolor{convbest}{rgb}{0.80,0.87,0.96}
\begin{tabular}{lccccc}
\toprule
Selection & 1\% & 5\% & 10\% & 20\% & 30\% \\
\midrule
High conviction & \cellcolor{convworst}$1.65 \pm 0.96$ & \cellcolor{convworst}$2.95 \pm 1.10$ & \cellcolor{convworst}$4.02 \pm 1.16$ & \cellcolor{convworst}$5.64 \pm 1.49$ & \cellcolor{convworst}$7.08 \pm 2.04$ \\
Low conviction & \cellcolor{convbest}$0.01 \pm 0.01$ & \cellcolor{convbest}$0.07 \pm 0.05$ & \cellcolor{convbest}$0.16 \pm 0.11$ & \cellcolor{convbest}$0.42 \pm 0.28$ & \cellcolor{convbest}$0.81 \pm 0.57$ \\
High degree & $1.36 \pm 0.77$ & $2.68 \pm 1.16$ & $3.48 \pm 1.24$ & $4.82 \pm 1.38$ & $5.77 \pm 1.57$ \\
Low degree & $0.04 \pm 0.03$ & $0.24 \pm 0.14$ & $0.54 \pm 0.30$ & $1.13 \pm 0.67$ & $1.88 \pm 1.10$ \\
Degree-matched (high) & $0.20 \pm 0.05$ & $0.52 \pm 0.21$ & $0.89 \pm 0.41$ & $1.29 \pm 0.76$ & $1.73 \pm 1.09$ \\
Degree-matched (low) & $0.07 \pm 0.04$ & $0.37 \pm 0.20$ & $0.67 \pm 0.41$ & $1.72 \pm 0.96$ & $2.92 \pm 1.52$ \\
High confidence & $0.33 \pm 0.13$ & $0.87 \pm 0.22$ & $1.37 \pm 0.33$ & $2.49 \pm 0.64$ & $3.52 \pm 1.04$ \\
Low confidence & $0.14 \pm 0.09$ & $0.48 \pm 0.21$ & $0.90 \pm 0.36$ & $1.90 \pm 0.75$ & $2.97 \pm 1.23$ \\
Random & $0.11 \pm 0.04$ & $0.43 \pm 0.20$ & $0.78 \pm 0.38$ & $1.69 \pm 0.83$ & $2.79 \pm 1.24$ \\
\bottomrule
\end{tabular}
\caption{Relative drop (\%) in the headline metric, accuracy or ROC-AUC on Minesweeper, Tolokers and Questions, expressed as a percentage of the baseline's headroom above chance, so that the two metrics are comparable. Intervention: pruning, as a function of the fraction of the training pool perturbed. Conviction score \texttt{last}. Each of the 8 datasets contributes one value, itself averaged over its folds. Mean $\pm$ standard error of the cross-dataset mean, the quoted spread being between-dataset variability. Larger is more damaging. In each column the most damaging selection is shaded red and the least damaging blue. }
\label{tab:conv_fwd_pruning_last}
\end{table}


% ------------------------------------------------------------------------------
\subsection{Complete decile tables}
\label{sec:supp_decile}

Tables~\ref{tab:conv_dec_masking-mean_last}--\ref{tab:conv_dec_pruning_last}
give the per-bin relative drop for the three remaining perturbation modes;
masking with zero replacement appears as
Table~\ref{tab:conv_dec_masking-zero_last} in the main text. The qualitative
pattern is the same throughout, and it is the one the experiment was built to
test: in all four modes the conviction column has both the smallest bin $1$ and
the largest bin $10$ of the three criteria, its lower bins are flat and close to
zero, and the bulk of its growth is concentrated in the last two bins. Neither
confidence nor degree produces a profile of that shape.

\input{sma_projet_files/Sections/conviction_tables/decile_masking-mean_GeneralSheaf_norm-true_conv-last.tex}

% auto-generated by quick_analysis/conviction_summary_tables.py
\begin{table}[H]
\centering
\small
\begin{tabular}{lccc}
\toprule
Bin & Conviction & Confidence & Degree \\
\midrule
1 & $0.18 \pm 0.09$ & $0.63 \pm 0.29$ & $0.57 \pm 0.27$ \\
2 & $0.22 \pm 0.14$ & $0.61 \pm 0.30$ & $0.59 \pm 0.29$ \\
3 & $0.25 \pm 0.24$ & $0.61 \pm 0.37$ & $0.56 \pm 0.34$ \\
4 & $0.39 \pm 0.32$ & $0.60 \pm 0.35$ & $0.54 \pm 0.33$ \\
5 & $0.59 \pm 0.35$ & $0.77 \pm 0.46$ & $0.51 \pm 0.35$ \\
6 & $0.63 \pm 0.38$ & $0.76 \pm 0.49$ & $0.52 \pm 0.36$ \\
7 & $0.70 \pm 0.41$ & $0.70 \pm 0.38$ & $0.66 \pm 0.32$ \\
8 & $0.87 \pm 0.45$ & $0.70 \pm 0.35$ & $0.56 \pm 0.32$ \\
9 & $1.07 \pm 0.51$ & $0.79 \pm 0.32$ & $0.75 \pm 0.37$ \\
10 & $2.77 \pm 0.66$ & $1.22 \pm 0.29$ & $1.92 \pm 0.45$ \\
\bottomrule
\end{tabular}
\caption{Relative drop (\%) in the headline metric, accuracy or ROC-AUC on Minesweeper, Tolokers and Questions, expressed as a percentage of the baseline's headroom above chance, so that the two metrics are comparable. Decile analysis under gating (hard): a single bin of the training pool is perturbed, with the pool split into 10 equal bins ranked by each criterion (bin 1 = lowest score, bin 10 = highest). Conviction score \texttt{last}. Mean $\pm$ standard error of the cross-dataset mean, over 8 datasets.}
\label{tab:conv_dec_gating-hard_last}
\end{table}


% auto-generated by quick_analysis/conviction_summary_tables.py
\begin{table}[H]
\centering
\small
\begin{tabular}{lccc}
\toprule
Bin & Conviction & Confidence & Degree \\
\midrule
1 & $0.16 \pm 0.11$ & $0.90 \pm 0.37$ & $0.54 \pm 0.30$ \\
2 & $0.23 \pm 0.16$ & $0.88 \pm 0.36$ & $0.55 \pm 0.34$ \\
3 & $0.30 \pm 0.27$ & $0.85 \pm 0.45$ & $0.60 \pm 0.38$ \\
4 & $0.51 \pm 0.35$ & $0.83 \pm 0.43$ & $0.57 \pm 0.38$ \\
5 & $0.69 \pm 0.42$ & $0.96 \pm 0.50$ & $0.66 \pm 0.41$ \\
6 & $0.74 \pm 0.45$ & $0.94 \pm 0.56$ & $0.65 \pm 0.44$ \\
7 & $0.85 \pm 0.48$ & $0.85 \pm 0.42$ & $0.73 \pm 0.39$ \\
8 & $1.00 \pm 0.52$ & $0.78 \pm 0.43$ & $0.68 \pm 0.38$ \\
9 & $1.23 \pm 0.59$ & $0.87 \pm 0.33$ & $0.98 \pm 0.46$ \\
10 & $4.02 \pm 1.16$ & $1.37 \pm 0.33$ & $3.52 \pm 1.25$ \\
\bottomrule
\end{tabular}
\caption{Relative drop (\%) in the headline metric, accuracy or ROC-AUC on Minesweeper, Tolokers and Questions, expressed as a percentage of the baseline's headroom above chance, so that the two metrics are comparable. Decile analysis under pruning: a single bin of the training pool is perturbed, with the pool split into 10 equal bins ranked by each criterion (bin 1 = lowest score, bin 10 = highest). Conviction score \texttt{last}. Mean $\pm$ standard error of the cross-dataset mean, over 8 datasets.}
\label{tab:conv_dec_pruning_last}
\end{table}


% ------------------------------------------------------------------------------
\subsection{Robustness under a different conviction score}
\label{sec:supp_robust}

The main text ranks nodes by \texttt{last}, the conviction at the final
diffusion layer, because that is the quantity the equilibrium theory speaks
about. A reader may reasonably worry that this rests on a single layer, whose
scale is arbitrary. We therefore repeat the whole analysis with \texttt{avg\_z},
in which the conviction of each post-diffusion layer $1,\dots,L$ is z-scored
separately and the results averaged, so that no single layer can dominate the
ranking. Tables~\ref{tab:conv_fwd_masking-zero_avg_z}
and~\ref{tab:conv_dec_masking-zero_avg_z} give the zero-masking sweep and decile
profile under this score, to be read against
Tables~\ref{tab:conv_fwd_masking-zero_last} and
\ref{tab:conv_dec_masking-zero_last} of the main text.

Every conclusion is unchanged, and the effect is if anything slightly sharper:
high conviction remains the most damaging selection and low conviction the least
at every budget, and the decile profile keeps the same shape. We report
\texttt{last} in the main text on theoretical grounds rather than because it is
the more favourable of the two. We also note an internal consistency check that
this comparison provides for free: the degree- and confidence-based rows are by
construction identical under the two scores, since they do not depend on the
conviction at all, and they do indeed come out identical.

% auto-generated by quick_analysis/conviction_summary_tables.py
\begin{table}[H]
\centering
\small
\definecolor{convworst}{rgb}{0.97,0.80,0.78}
\definecolor{convbest}{rgb}{0.80,0.87,0.96}
\begin{tabular}{lccccc}
\toprule
Selection & 1\% & 5\% & 10\% & 20\% & 30\% \\
\midrule
High conviction & \cellcolor{convworst}$0.56 \pm 0.17$ & \cellcolor{convworst}$1.68 \pm 0.57$ & \cellcolor{convworst}$2.65 \pm 0.98$ & \cellcolor{convworst}$4.10 \pm 1.61$ & \cellcolor{convworst}$5.18 \pm 2.21$ \\
Low conviction & \cellcolor{convbest}$0.01 \pm 0.01$ & \cellcolor{convbest}$0.02 \pm 0.03$ & \cellcolor{convbest}$0.02 \pm 0.06$ & \cellcolor{convbest}$0.13 \pm 0.18$ & \cellcolor{convbest}$0.34 \pm 0.29$ \\
High degree & $0.37 \pm 0.16$ & $0.75 \pm 0.28$ & $1.06 \pm 0.40$ & $1.76 \pm 0.69$ & $2.46 \pm 0.98$ \\
Low degree & $0.03 \pm 0.02$ & $0.19 \pm 0.10$ & $0.38 \pm 0.21$ & $0.85 \pm 0.44$ & $1.35 \pm 0.74$ \\
Degree-matched (high) & $0.07 \pm 0.04$ & $0.19 \pm 0.15$ & $0.27 \pm 0.26$ & $0.52 \pm 0.49$ & $0.77 \pm 0.55$ \\
Degree-matched (low) & $0.06 \pm 0.03$ & $0.26 \pm 0.12$ & $0.49 \pm 0.33$ & $1.17 \pm 0.66$ & $2.22 \pm 1.23$ \\
High confidence & $0.15 \pm 0.07$ & $0.55 \pm 0.16$ & $0.92 \pm 0.28$ & $1.66 \pm 0.54$ & $2.50 \pm 0.89$ \\
Low confidence & $0.04 \pm 0.03$ & $0.17 \pm 0.12$ & $0.35 \pm 0.21$ & $0.80 \pm 0.43$ & $1.32 \pm 0.69$ \\
Random & $0.07 \pm 0.04$ & $0.21 \pm 0.15$ & $0.43 \pm 0.29$ & $1.04 \pm 0.55$ & $1.69 \pm 0.89$ \\
\bottomrule
\end{tabular}
\caption{Relative drop (\%) in the headline metric, accuracy or ROC-AUC on Minesweeper, Tolokers and Questions, expressed as a percentage of the baseline's headroom above chance, so that the two metrics are comparable. Intervention: masking (zero), as a function of the fraction of the training pool perturbed. Conviction score \texttt{avg\_z}. Each of the 8 datasets contributes one value, itself averaged over its folds. Mean $\pm$ standard error of the cross-dataset mean, the quoted spread being between-dataset variability. Larger is more damaging. In each column the most damaging selection is shaded red and the least damaging blue. }
\label{tab:conv_fwd_masking-zero_avg_z}
\end{table}


% auto-generated by quick_analysis/conviction_summary_tables.py
\begin{table}[H]
\centering
\small
\begin{tabular}{lccc}
\toprule
Bin & Conviction & Confidence & Degree \\
\midrule
1 & $0.02 \pm 0.06$ & $0.35 \pm 0.21$ & $0.38 \pm 0.21$ \\
2 & $0.05 \pm 0.12$ & $0.36 \pm 0.21$ & $0.42 \pm 0.22$ \\
3 & $0.12 \pm 0.12$ & $0.39 \pm 0.26$ & $0.41 \pm 0.24$ \\
4 & $0.16 \pm 0.19$ & $0.42 \pm 0.25$ & $0.37 \pm 0.27$ \\
5 & $0.23 \pm 0.23$ & $0.48 \pm 0.38$ & $0.45 \pm 0.29$ \\
6 & $0.35 \pm 0.24$ & $0.58 \pm 0.41$ & $0.46 \pm 0.31$ \\
7 & $0.41 \pm 0.37$ & $0.61 \pm 0.37$ & $0.55 \pm 0.29$ \\
8 & $0.62 \pm 0.37$ & $0.58 \pm 0.37$ & $0.44 \pm 0.31$ \\
9 & $1.06 \pm 0.50$ & $0.61 \pm 0.29$ & $0.56 \pm 0.34$ \\
10 & $2.65 \pm 0.98$ & $0.92 \pm 0.28$ & $1.06 \pm 0.40$ \\
\bottomrule
\end{tabular}
\caption{Relative drop (\%) in the headline metric, accuracy or ROC-AUC on Minesweeper, Tolokers and Questions, expressed as a percentage of the baseline's headroom above chance, so that the two metrics are comparable. Decile analysis under masking (zero): a single bin of the training pool is perturbed, with the pool split into 10 equal bins ranked by each criterion (bin 1 = lowest score, bin 10 = highest). Conviction score \texttt{avg\_z}. Mean $\pm$ standard error of the cross-dataset mean, over 8 datasets.}
\label{tab:conv_dec_masking-zero_avg_z}
\end{table}


% ------------------------------------------------------------------------------
\subsection{Graph integrity under pruning}
\label{sec:supp_prune_diag}

Pruning is the only intervention that changes the graph, so for every
(strategy, budget) pair we also record the fraction of nodes and edges
remaining, the number of connected components, the size of the giant component
and the class balance within each split, before and after. This lets the
performance drop of Table~\ref{tab:conv_fwd_pruning_last} be read against the
structural damage that produced it, which is what
Table~\ref{tab:conv_prune_edges_last} reports.

The comparison matters most for the degree-based strategies, and we would rather
verify the point than assert it. Removing the top-$k\%$ by degree deletes by
construction the largest possible number of edges: reading
Table~\ref{tab:conv_prune_edges_last}, $51.6 \pm 4.3\%$ of the graph at the
$30\%$ budget, against $39.4 \pm 2.8\%$ for the top-$k\%$ by conviction.
The two criteria therefore do not pay the same structural price for the same
budget, and Table~\ref{tab:conv_prune_evd_last} sets the two prices side by side
with what each of them buys.

Read per dataset in Table~\ref{tab:conv_prune_evd_last}, the picture is
unambiguous. At the $30\%$ budget the degree
selection deletes a larger fraction of the edges than the conviction selection
on all eight datasets, and on four of them, Amazon Ratings, Citeseer,
Minesweeper and Roman Empire, it nonetheless does less damage. The clearest case
is Citeseer, where degree deletes $48.1\%$ of the edges for a drop of $1.01\%$
while conviction deletes $37.0\%$ for a drop of $3.27\%$: a third fewer edges,
three times the damage. Minesweeper shows the same thing with larger numbers,
$33.2\%$ of edges for a drop of $8.19\%$ against $29.8\%$ for $17.15\%$. On the
remaining four datasets the degree selection does come out ahead on the drop,
but it does so while cutting far more of the graph, most extremely on Tolokers,
where it deletes $65.2\%$ of the edges against $40.4\%$ for a drop of $3.17\%$
against $2.63\%$. Fragmentation tells the same story from the other side, and is
reported in Table~\ref{tab:conv_prune_comps_last}: at the $30\%$ budget the
degree selection leaves the graph in $3225$ connected components on average,
against $2142$ for the conviction selection.

We should be equally explicit about a limitation of the degree-matched control
that the same diagnostics expose. The control is built to reproduce the degree
distribution of the high-conviction group, and at small budgets it does so and
deletes a comparable share of the graph. Reading
Table~\ref{tab:conv_prune_edges_last} at the $1\%$ budget, it deletes
$5.0 \pm 1.3\%$ against $7.7 \pm 3.1\%$ for the high-conviction group and five
times what a random selection of the same size deletes, while costing, in
Table~\ref{tab:conv_fwd_pruning_last}, $0.20 \pm 0.05\%$ against
$1.65 \pm 0.96\%$. At the largest budgets the matching is diluted,
because the pool of unused high-degree candidates is exhausted, and by the
$30\%$ budget the control deletes $28.3 \pm 1.7\%$ of the edges, which is
indistinguishable from the $28.3 \pm 1.0\%$ deleted by a random selection. The
degree-matched comparison is therefore at its strongest at the small budgets,
and should be read with that in mind at the large ones.

% auto-generated by quick_analysis/conviction_summary_tables.py
\begin{table}[H]
\centering
\small
\begin{tabular}{lccccc}
\toprule
Selection & 1\% & 5\% & 10\% & 20\% & 30\% \\
\midrule
High conviction & $7.7 \pm 3.1$ & $15.9 \pm 3.5$ & $22.0 \pm 3.4$ & $31.3 \pm 3.1$ & $39.4 \pm 2.8$ \\
Low conviction & $0.5 \pm 0.1$ & $2.7 \pm 0.5$ & $5.6 \pm 1.0$ & $11.6 \pm 1.8$ & $18.0 \pm 2.5$ \\
High degree & $10.2 \pm 3.4$ & $23.5 \pm 5.1$ & $32.7 \pm 5.5$ & $44.2 \pm 5.1$ & $51.6 \pm 4.3$ \\
Low degree & $0.4 \pm 0.1$ & $1.9 \pm 0.6$ & $4.2 \pm 1.2$ & $8.7 \pm 2.5$ & $13.3 \pm 3.7$ \\
Degree-matched (high) & $5.0 \pm 1.3$ & $11.5 \pm 1.4$ & $16.6 \pm 1.3$ & $23.6 \pm 1.1$ & $28.3 \pm 1.7$ \\
Degree-matched (low) & $0.5 \pm 0.1$ & $2.8 \pm 0.6$ & $5.9 \pm 1.1$ & $12.4 \pm 1.9$ & $19.3 \pm 2.6$ \\
High confidence & $2.0 \pm 0.4$ & $6.9 \pm 1.1$ & $11.7 \pm 1.7$ & $20.3 \pm 2.6$ & $28.1 \pm 3.1$ \\
Low confidence & $1.7 \pm 0.9$ & $7.1 \pm 2.6$ & $12.2 \pm 3.5$ & $20.7 \pm 4.2$ & $28.2 \pm 4.5$ \\
Random & $1.0 \pm 0.1$ & $5.0 \pm 0.2$ & $10.0 \pm 0.4$ & $19.5 \pm 0.7$ & $28.3 \pm 1.0$ \\
\bottomrule
\end{tabular}
\caption{Structural cost of pruning: percentage of the graph's edges deleted together with the selected nodes, as a function of the fraction of the training pool removed. Mean $\pm$ standard error over the 8 datasets. To be read against Table~\ref{tab:conv_fwd_pruning_last}, which gives what each of these selections costs in performance: a strategy that deletes more edges for less damage is destroying structure rather than information.}
\label{tab:conv_prune_edges_last}
\end{table}


% auto-generated by quick_analysis/conviction_summary_tables.py
\begin{table}[H]
\centering
\small
\begin{tabular}{lccccc}
\toprule
Selection & 1\% & 5\% & 10\% & 20\% & 30\% \\
\midrule
High conviction & $1167$ & $1596$ & $1742$ & $1960$ & $2142$ \\
Low conviction & $74$ & $106$ & $144$ & $224$ & $322$ \\
High degree & $1057$ & $1859$ & $2373$ & $2959$ & $3225$ \\
Low degree & $66$ & $75$ & $87$ & $112$ & $133$ \\
Degree-matched (high) & $248$ & $361$ & $487$ & $620$ & $693$ \\
Degree-matched (low) & $71$ & $104$ & $151$ & $259$ & $386$ \\
High confidence & $184$ & $270$ & $342$ & $504$ & $687$ \\
Low confidence & $160$ & $464$ & $749$ & $1151$ & $1421$ \\
Random & $92$ & $212$ & $351$ & $646$ & $906$ \\
\bottomrule
\end{tabular}
\caption{Fragmentation under pruning: number of connected components left in the graph after the selected nodes are deleted, averaged over the ten folds and then over the 8 datasets. The companion of Table~\ref{tab:conv_prune_edges_last}: it measures the same structural damage from the other side, since a selection may delete few edges and still cut the graph into many pieces.}
\label{tab:conv_prune_comps_last}
\end{table}


% auto-generated by quick_analysis/conviction_summary_tables.py
\begin{table}[H]
\centering
\small
\begin{tabular}{lcccc}
\toprule
 & \multicolumn{2}{c}{High conviction} & \multicolumn{2}{c}{High degree} \\
\cmidrule(lr){2-3}\cmidrule(lr){4-5}
Dataset & Edges & Drop & Edges & Drop \\
\midrule
Amazon Ratings & $32.9$ & $6.43$ & $46.1$ & $5.88$ \\
Citeseer & $37.0$ & $3.27$ & $48.1$ & $1.01$ \\
Cora & $36.3$ & $2.58$ & $46.1$ & $2.68$ \\
Minesweeper & $29.8$ & $17.15$ & $33.2$ & $8.19$ \\
Pubmed & $49.6$ & $1.34$ & $61.5$ & $1.45$ \\
Questions & $52.4$ & $10.31$ & $68.3$ & $12.31$ \\
Roman Empire & $36.5$ & $12.92$ & $44.4$ & $11.48$ \\
Tolokers & $40.4$ & $2.63$ & $65.2$ & $3.17$ \\
\bottomrule
\end{tabular}
\caption{Structural cost against performance cost, per dataset, for pruning at the 30\% budget. ``Edges'' is the percentage of the graph's edges deleted along with the selected nodes; ``Drop'' is the relative drop (\%) in the headline metric, as everywhere else in this section. Selecting by degree deletes the larger fraction of the graph on every dataset, yet on several of them it costs less than selecting by conviction, which is what shows that part of the damage a degree criterion does is structural rather than informational. Conviction score \texttt{last}, values averaged over the ten folds.}
\label{tab:conv_prune_evd_last}
\end{table}


% ------------------------------------------------------------------------------
\subsection{Soft gating}
\label{sec:supp_soft}

Besides the three perturbation modes of Section~\ref{sec:conviction} we ran a
\emph{soft} gating variant, which we report here because it is a negative
result and no claim in the main text rests on it. Unlike every other experiment,
it has no budget and no selection axis: the whole training pool is attenuated at
once, each node's outgoing message in proportion to its own score, the
lowest-scoring node being suppressed the most. Table~\ref{tab:conv_soft_last}
gives the outcome. Attenuating by conviction costs $3.09 \pm 0.95\%$, by degree
$4.19 \pm 1.95\%$ and by confidence $0.67 \pm 0.37\%$, with error bars so wide
that the three cannot be separated. This is expected in hindsight, since the
soft variant simultaneously suppresses the low-conviction nodes and preserves
the high-conviction ones, so the two effects that the budget sweep isolates are
here superimposed and partly cancel. We include the table so that the negative
result is as inspectable as the positive ones.

% auto-generated by quick_analysis/conviction_summary_tables.py
\begin{table}[H]
\centering
\small
\begin{tabular}{lc}
\toprule
Indicator & Relative drop (\%) \\
\midrule
Confidence & $0.67 \pm 0.37$ \\
Conviction & $3.09 \pm 0.95$ \\
Degree & $4.19 \pm 1.95$ \\
\bottomrule
\end{tabular}
\caption{Relative drop (\%) in the headline metric, accuracy or ROC-AUC on Minesweeper, Tolokers and Questions, expressed as a percentage of the baseline's headroom above chance, so that the two metrics are comparable. Soft gating: instead of selecting a subset, every training node's outgoing message is attenuated in proportion to its own score, the lowest-scoring node being suppressed the most (gate floor $g_{\min}=0.1$). There is therefore no budget axis and one number per indicator. Conviction score \texttt{last}. Mean $\pm$ standard error of the cross-dataset mean, over 8 datasets.}
\label{tab:conv_soft_last}
\end{table}

